\chapter{Lezione 7 --- Anti-trasformata z, fratti semplici e funzione di trasferimento discreta}

\section{Richiamo: trasformata z e sistemi discreti}

Riprendiamo la definizione di trasformata z monolatera di una successione
\[
x(k), \qquad k=0,1,2,\dots
\]
data da
\[
X(z)=\mathcal{Z}\{x(k)\}=\sum_{k=0}^{+\infty} x(k)\,z^{-k}.
\]

Nel caso di un segnale continuo \(f(t)\) campionato con periodo \(T\), si scrive pi\`u rigorosamente
\[
\mathcal{Z}\{f(kT)\}=\sum_{k=0}^{+\infty} f(kT)\,z^{-k}.
\]

Negli appunti si sottolinea che la trasformata z nasce come riscrittura della trasformata di Laplace
del segnale campionato, con il cambio di variabile
\[
z=e^{sT}.
\]

\section{Richiamo: matrice e funzione di trasferimento}

Per un sistema discreto in forma di stato
\[
\begin{cases}
x(k+1)=Ax(k)+Bu(k),\\[4pt]
y(k)=Cx(k),
\end{cases}
\]
la matrice di trasferimento \`e
\[
G(z)=C(zI-A)^{-1}B.
\]

Nel caso \textbf{SISO}, si ha
\[
u(k)\in\mathbb{R}, \qquad y(k)\in\mathbb{R}, \qquad x(k)\in\mathbb{R}^n,
\]
e quindi
\[
G(z)=c^T(zI-A)^{-1}b=\frac{N(z)}{D(z)},
\]
cio\`e una funzione razionale in \(z\).

In condizioni iniziali nulle vale
\[
Y(z)=G(z)U(z).
\]

Il problema affrontato in questa lezione \`e allora il seguente:
dato \(Y(z)\), oppure pi\`u in generale data una funzione razionale \(G(z)\),
\emph{come si torna alla successione nel tempo discreto}?

\section{Questione di unicit\`a}

Negli appunti compare un'osservazione importante sul confronto tra Laplace e zeta.

Nel continuo, due funzioni che differiscono solo su un insieme di misura nulla
possono avere la stessa trasformata di Laplace, perch\'e la trasformata si basa su un integrale.

Nel discreto, invece, la successione \`e definita punto per punto:
se due successioni sono diverse, differiscono necessariamente su almeno un campione,
e quindi in condizioni standard danno trasformate z diverse.

Per questo, in ambito discreto, si ragiona di fatto con una corrispondenza biunivoca
tra successione e trasformata z, purch\'e si tenga conto anche della regione di convergenza.

\section{Metodi per l'anti-trasformata z}

Negli appunti vengono richiamati due metodi principali per ricavare la successione \(x(k)\)
a partire da \(X(z)\), quando \(X(z)\) \`e razionale:

\begin{enumerate}
    \item \textbf{fratti semplici};
    \item \textbf{lunga divisione}.
\end{enumerate}

Il primo \`e il metodo principale quando il denominatore si pu\`o fattorizzare bene
e si vogliono ricavare formule chiuse.
Il secondo \`e utile quando si vuole sviluppare \(X(z)\) come serie in potenze di \(z^{-1}\)
e leggere direttamente i coefficienti della successione.

\section{Richiamo sui fratti semplici nel dominio di Laplace}

Negli appunti si fa prima un parallelo con la trasformata di Laplace.

Se
\[
G(s)=\frac{N(s)}{D(s)}
\]
e il denominatore si fattorizza come
\[
D(s)=\prod_{i=1}^{n}(s-\alpha_i),
\]
con radici semplici, allora si pu\`o scrivere
\[
G(s)=\sum_{i=1}^{n}\frac{A_i}{s-\alpha_i}.
\]

Da qui segue immediatamente
\[
g(t)=\sum_{i=1}^{n} A_i e^{\alpha_i t}.
\]

Se invece una radice \(\alpha_i\) ha molteplicit\`a maggiore di \(1\), compaiono termini del tipo
\[
\frac{A_{i,1}}{s-\alpha_i}+\frac{A_{i,2}}{(s-\alpha_i)^2}+\dots+\frac{A_{i,\mu_i}}{(s-\alpha_i)^{\mu_i}}.
\]

I coefficienti si possono determinare:
\begin{itemize}
    \item per confronto dei coefficienti;
    \item con il metodo dei residui, nel caso di poli semplici:
    \[
    A_i=\lim_{s\to \alpha_i}(s-\alpha_i)G(s).
    \]
\end{itemize}

\section{Applicazione del metodo dei fratti semplici nel dominio \(z\)}

L'idea \`e del tutto analoga nel dominio z.

Se
\[
G(z)=\frac{N(z)}{D(z)},
\]
si vorrebbe decomporre \(G(z)\) in termini elementari di cui si conosce gi\`a
l'anti-trasformata z.

\subsection*{Osservazione pratica}

Negli appunti si osserva che spesso, scrivendo \(G(z)\) in funzione di \(z\),
il numeratore e il denominatore hanno lo stesso grado.
In questi casi, per applicare comodamente i fratti semplici,
si lavora su
\[
\frac{G(z)}{z}=\frac{N(z)}{z\,D(z)},
\]
cio\`e si aggiunge un polo nell'origine.

Questo rende pi\`u naturale la decomposizione in termini elementari
ed \`e coerente con il fatto che molte tavole di anti-trasformata z
sono espresse in forme del tipo
\[
\frac{z}{z-a}, \qquad
\frac{z\sin\theta}{z^2-2\cos\theta\,z+1},
\qquad
\frac{z(z-\cos\theta)}{z^2-2\cos\theta\,z+1}.
\]

\section{Tavola di riferimento per l'anti-trasformata z}

Le forme elementari da tenere a mente sono:

\[
\mathcal{Z}\{\delta(k)\}=1,
\]
\[
\mathcal{Z}\{a^k 1(k)\}=\frac{z}{z-a},
\qquad |z|>|a|,
\]
\[
\mathcal{Z}\{\sin(k\theta)\,1(k)\}
=
\frac{z\sin\theta}{z^2-2\cos\theta\,z+1},
\]
\[
\mathcal{Z}\{\cos(k\theta)\,1(k)\}
=
\frac{z(z-\cos\theta)}{z^2-2\cos\theta\,z+1}.
\]

Queste forme sono quelle che permettono di riconoscere rapidamente
i contributi elementari nella decomposizione in fratti semplici.

\section{Esempio di decomposizione nel dominio \(z\)}

Negli appunti compare l'esempio
\[
G_1(z)=\frac{z+1}{\left(z-\frac13\right)\left(z^2-\frac14 z+1\right)}.
\]

I poli sono:
\[
z_1=\frac13,
\qquad
z_{2,3}=\frac18 \pm j\,\frac{3\sqrt7}{8}.
\]

Poich\'e
\[
z^2-\frac14 z+1
=
z^2-2\cos\theta\,z+1,
\]
si ricava
\[
2\cos\theta=\frac14
\qquad\Longrightarrow\qquad
\cos\theta=\frac18,
\qquad
\theta=\arccos\!\left(\frac18\right).
\]

Per applicare il metodo, si considera
\[
\frac{G_1(z)}{z}
=
\frac{z+1}{z\left(z-\frac13\right)\left(z^2-\frac14 z+1\right)}.
\]

Una forma di decomposizione coerente con le tavole note \`e
\[
\frac{G_1(z)}{z}
=
\frac{A}{z}
+
\frac{B}{z-\frac13}
+
C\,\frac{z\sin\theta}{z^2-\frac14 z+1}
+
D\,\frac{z(z-\cos\theta)}{z^2-\frac14 z+1}.
\]

I coefficienti \(A,B,C,D\) si determinano per uguaglianza tra i due membri,
oppure combinando confronto dei coefficienti e uso delle forme note.

Una volta trovati i coefficienti, l'anti-trasformata si legge termine a termine come combinazione di:
\begin{itemize}
    \item un termine impulsivo;
    \item un'esponenziale discreta \(\left(\frac13\right)^k\);
    \item un termine sinusoidale \(\sin(k\theta)\);
    \item un termine cosinusoidale \(\cos(k\theta)\).
\end{itemize}

\subsection*{Osservazione}

Negli appunti \`e evidenziato che i vari metodi possono anche essere \emph{combinati}:
per esempio si possono isolare alcuni termini con i fratti semplici
e ricavare altri sviluppando in serie.

\section{Metodo della lunga divisione}

Il secondo metodo consiste nello sviluppare \(X(z)\) in serie di potenze di \(z^{-1}\):
\[
X(z)=x(0)+x(1)z^{-1}+x(2)z^{-2}+\dots
\]

Se si riesce a scrivere la funzione razionale in questa forma,
allora i coefficienti della serie sono direttamente i valori della successione.

Questo metodo \`e particolarmente utile:
\begin{itemize}
    \item quando si vogliono calcolare i primi campioni;
    \item quando la fattorizzazione del denominatore \`e scomoda;
    \item quando interessa una realizzazione ricorsiva.
\end{itemize}

\section{Collegamento con la funzione di trasferimento discreta}

Nel caso SISO, la funzione di trasferimento
\[
G(z)=\frac{N(z)}{D(z)}
\]
descrive il legame tra ingresso e uscita:
\[
Y(z)=G(z)U(z).
\]

Per tornare al dominio discreto si possono seguire due strade:
\begin{enumerate}
    \item anti-trasformare direttamente \(Y(z)\) con fratti semplici o lunga divisione;
    \item riscrivere \(G(z)\) in forma polinomiale in \(z^{-1}\) e ottenere
    l'equazione alle differenze.
\end{enumerate}

La seconda strada \`e quella pi\`u utile quando si vuole implementare il sistema,
mentre la prima \`e quella pi\`u utile quando si vuole trovare una formula esplicita
per la risposta.

\section{Osservazione sul continuo e sul discreto}

Negli appunti compare un'ultima osservazione concettuale:
nel continuo esistono formalmente funzioni diverse che differiscono solo su insiemi
di misura nulla e che quindi hanno la stessa trasformata integrale.

Nel discreto questo problema, dal punto di vista pratico, non si presenta:
i segnali reali sono successioni di valori campionati, quindi l'informazione
\`e contenuta esattamente nei campioni.

Per questo, nell'uso ingegneristico della trasformata z,
si ragiona come se ci fosse una corrispondenza biunivoca tra sequenza e trasformata.

\section{Sintesi della lezione}

I punti principali della lezione sono:
\begin{enumerate}
    \item richiamo della matrice di trasferimento discreta
    \[
    G(z)=C(zI-A)^{-1}B;
    \]
    \item nel caso SISO, \(G(z)\) \`e una funzione razionale
    \[
    G(z)=\frac{N(z)}{D(z)};
    \]
    \item il problema dell'anti-trasformata z consiste nel tornare da \(X(z)\)
    alla successione \(x(k)\);
    \item i due metodi principali sono:
    \begin{itemize}
        \item fratti semplici;
        \item lunga divisione;
    \end{itemize}
    \item il metodo dei fratti semplici nel dominio \(z\) ricalca quello della Laplace;
    \item spesso conviene lavorare su
    \[
    \frac{G(z)}{z}
    \]
    per ottenere una forma pi\`u comoda da decomporre;
    \item i termini elementari dell'anti-trasformata sono esponenziali discreti,
    seni e coseni discreti, oltre agli impulsi;
    \item l'esempio svolto mostra come trattare un polo reale e una coppia di poli
    complessi coniugati nel piano \(z\).
\end{enumerate}